\documentclass[11pt]{article}

\usepackage[T1]{fontenc}
\usepackage[margin=1in]{geometry}
\usepackage{amsmath,amssymb}
\usepackage[hidelinks]{hyperref}

% Working draft: replace this provisional title and add author information later.
% This is a standalone draft, not a conference-formatted submission.
\title{Hierarchical Learning and Control for Synchronized Arrival\newline
under a Shared Deadline}
\author{}
\date{}

\begin{document}
\maketitle

\begin{abstract}
We present a hierarchical learning and control framework for synchronized
arrival of mobile robots at arbitrary preassigned goals under a shared
deadline. A shared policy, trained using Multi-Agent Proximal Policy
Optimization (MAPPO) with centralized training and decentralized execution,
generates local waypoints and preferred speeds from each vehicle's own state,
goal, timing information, and locally sensed neighbours. Pure-pursuit rollout
converts these decisions into timed reference trajectories, which independent
nonlinear model predictive control (NMPC) controllers track while accounting
for actuator limits and predicted inter-vehicle separation. The controller
operates inside the training loop, allowing the policy to learn coordination
under the constraints and interventions encountered during execution. A
masked attention encoder aggregates the nearest neighbours into a fixed-size
representation, allowing the policy architecture to accommodate varying
neighbourhood sizes. Scalability and robustness are assessed empirically.
Online decisions require neither a centralized planner nor exchanged
neighbour goals or future trajectories, given the supplied common deadline.
We assess the contributions of learning and constrained control through
matched comparisons of pure-pursuit execution, MPC added after policy
training, and training with MPC in the loop, alongside a non-learning deadline
controller. These experiments characterize the trade-offs between synchronized
arrival, collision avoidance, and completion failures.
\end{abstract}

\section{Introduction and Related Work}
\label{sec:introduction}
% Add the introduction, related work, and verified citations here later.

\section{Contribution}
\label{sec:contribution}

The contribution is learning decentralized motion references that coordinate
both route selection and speed so that mobile robots reach different assigned
goals at a common deadline, while accounting for the behaviour of a
constrained low-level controller.

The architecture uses a hierarchical combination of RL and NMPC to address
the problem of simultaneous arrival, with robustness and scalability
evaluated experimentally.

\section{Problem Formulation and Background}
\label{sec:background}

Multi-agent reinforcement learning (MARL) is based on reinforcement learning
(RL), in which an agent interacts with its environment and learns a policy
that aims to maximize the expected cumulative reward.

The problem is modelled as a decentralized partially observable Markov
decision process (Dec-POMDP). In this work, with homogeneous agents and
deterministic local observations, we write the Dec-POMDP as
\begin{equation}
    \mathcal{M}
    = (\mathcal{S},\mathcal{A},\mathcal{O},R,P,\Omega,n,\gamma),
\end{equation}
where $\mathcal{S}$ is the global state space, $\mathcal{A}$ is the common
action space of each agent, $\mathcal{O}$ is the local observation space,
and $n$ is the number of agents. The local observation of agent $i$ at time
step $t$ is
\begin{equation}
    o_{i,t} = \Omega(s_t,i).
\end{equation}
The joint action is
\begin{equation}
    \mathbf{a}_t = (a_{1,t},\ldots,a_{n,t}) \in \mathcal{A}^{n}.
\end{equation}
The transition model $P(s_{t+1}\mid s_t,\mathbf{a}_t)$ describes the
probability of transitioning from $s_t$ to $s_{t+1}$ under the joint action.
The shared reward is $r_t=R(s_t,\mathbf{a}_t)$, and $\gamma\in[0,1)$ is
the discount factor.

Agents use a shared policy $\pi_\theta(a_{i,t}\mid o_{i,t})$, parameterized
by $\theta$, to produce an action from their own local observation. They
jointly optimize the expected discounted cumulative reward:
\begin{equation}
    J(\theta)
    = \mathbb{E}_{\pi_\theta,P}\!\left[
        \sum_{t=0}^{H_{\mathrm{ep}}-1}\gamma^t r_t
      \right],
    \label{eq:rl-objective}
\end{equation}
where $H_{\mathrm{ep}}$ is the number of steps in an episode.

MARL algorithms can be categorized according to the information available
during training and execution. During training, an algorithm may be
restricted to the local information observed by each agent
(decentralized training), or it may use information about all agents in
the multi-agent system (centralized training).

Centralized training with decentralized execution (CTDE) combines these
information settings: training can use global information, while each agent
uses only its permitted local information to make decisions during execution.

MAPPO is an on-policy deep reinforcement learning algorithm that adopts
the CTDE paradigm. In our implementation, a shared actor
$\pi_\theta(a_{i,t}\mid o_{i,t})$ receives only agent $i$'s local
observation and outputs that agent's action. A centralized critic with
parameters $\phi$ receives the global state $s_t$, which is not supplied
to the actor during execution. The critic is used during training to
estimate values and compute advantages using generalized advantage
estimation (GAE). We write its output as $V_\phi(s_t,i)$ because the
implementation produces one value estimate per agent from the global state.

Initially, the algorithm collects a sequence of experiences called a
rollout. It stores local observations, global states, sampled latent actions,
old action log-probabilities, critic values, rewards, next-state values,
termination and rollout-boundary flags, and active-agent masks. The latent
Gaussian actions are transformed using $\tanh$ to obtain bounded policy
actions. The policy log-probabilities account for this transformation.

PPO needs to measure how much the new policy differs from the policy that
generated the rollout. It therefore uses the probability ratio
\begin{align}
    \rho_{i,t}(\theta)
    &= \frac{\pi_\theta(a_{i,t}\mid o_{i,t})}
            {\pi_{\theta_{\mathrm{old}}}(a_{i,t}\mid o_{i,t})}
       \label{eq:probability-ratio} \\
    &= \exp\!\left(
        \log\pi_\theta(a_{i,t}\mid o_{i,t})
        -\log\pi_{\theta_{\mathrm{old}}}(a_{i,t}\mid o_{i,t})
       \right),
\end{align}
where $\theta_{\mathrm{old}}$ denotes the parameters used to collect the
rollout. Both probabilities are evaluated for the same collected action
and observation.

The discounted return from time step $t$ is
\begin{equation}
    G_t
    = \sum_{k=t}^{H_{\mathrm{ep}}-1}\gamma^{k-t}r_k
    = r_t + \gamma r_{t+1} + \gamma^2 r_{t+2} + \cdots.
    \label{eq:discounted-return}
\end{equation}
For the collected rollout, GAE estimates how much better an action was
than the critic's value prediction. First, we compute the temporal-difference
error:
\begin{equation}
    \delta_{i,t}
    = r_t
      + \gamma(1-d_t)V_{\phi_{\mathrm{old}}}(s_{t+1},i)
      - V_{\phi_{\mathrm{old}}}(s_t,i),
    \label{eq:td-error}
\end{equation}
where $d_t=1$ if the task terminates after the transition, and $d_t=0$
otherwise. Thus, a true terminal transition does not bootstrap from a
next-state value. The subscript $\phi_{\mathrm{old}}$ indicates that
these are the fixed critic estimates collected before the update.

The GAE advantages are computed backwards through the rollout:
\begin{equation}
    \widehat{A}_{i,t}
    = \delta_{i,t}
      + \gamma\lambda(1-d_t)(1-b_t)\widehat{A}_{i,t+1},
    \qquad \widehat{A}_{i,H}=0,
    \label{eq:gae}
\end{equation}
where $H$ is the rollout length, $\lambda\in[0,1]$ controls the GAE
weighting, and $b_t=1$ marks an artificial rollout or curriculum boundary.
Such a boundary stops the advantage recursion but retains the next-state
bootstrap in~\eqref{eq:td-error} when the task has not terminated.

The critic's regression target is the GAE-based return estimate
\begin{equation}
    \widehat{G}_{i,t}
    = \widehat{A}_{i,t} + V_{\phi_{\mathrm{old}}}(s_t,i).
    \label{eq:value-target}
\end{equation}
The critic is trained by minimizing the value loss
\begin{equation}
    \mathcal{L}_{V}(\phi)
    = \frac{1}{2}\,
      \mathbb{E}_{(i,t)\sim\mathcal{B}}\!\left[
        \bigl(V_\phi(s_t,i)-\widehat{G}_{i,t}\bigr)^2
      \right],
    \label{eq:critic-loss}
\end{equation}
where $\mathcal{B}$ denotes the collected agent--time samples. The return
targets are held fixed during the optimization update.

The actor uses the advantages to form the PPO clipped objective. Let
$\widetilde{A}_{i,t}$ denote the advantages normalized over active decision
samples when their variance is nonzero. Then
\begin{equation}
    \mathcal{L}_{\mathrm{clip}}(\theta)
    = \mathbb{E}_{(i,t)\sim\mathcal{B}_{\mathrm{active}}}\!\left[
        \min\!\left(
          \rho_{i,t}(\theta)\widetilde{A}_{i,t},\,
          \operatorname{clip}\!\left(
            \rho_{i,t}(\theta),1-\epsilon,1+\epsilon
          \right)\widetilde{A}_{i,t}
        \right)
      \right],
    \label{eq:ppo-clipping}
\end{equation}
where $\epsilon>0$ is the clipping parameter and
$\mathcal{B}_{\mathrm{active}}$ contains samples from agents that were
active when their actions were selected. The actor maximizes this objective.
Clipping limits the incentive for large policy changes in the surrogate
objective; it does not impose a hard bound on every probability ratio.

The implementation combines the actor objective, critic loss, and an entropy
bonus into the loss that is minimized:
\begin{equation}
    \mathcal{L}(\theta,\phi)
    = -\mathcal{L}_{\mathrm{clip}}(\theta)
      + c_V\mathcal{L}_{V}(\phi)
      - c_H\,\mathbb{E}_{(i,t)\sim\mathcal{B}_{\mathrm{active}}}\!\left[
          \mathcal{H}\bigl(\pi_\theta(\cdot\mid o_{i,t})\bigr)
        \right],
    \label{eq:combined-loss}
\end{equation}
where $c_V$ and $c_H$ weight the value loss and policy entropy, respectively.
The entropy term encourages exploration. The actor and critic are updated
using minibatches from the rollout, while the stored old log-probabilities,
advantages, and return targets remain fixed during these updates.

\end{document}
